\section{Methods}
\label{sec:Methods}

We adopt the formalism of relativistic quantum field theory on a flat,
non-interacting Minkowski spacetime. The background geometry is defined by the
orthonormal basis $\{ \gamma^\mu \}$ satisfying the Clifford algebra
$\mathcal{C}\ell_{1,3}(\mathbb{C})$, given by the anticommutation relation
\begin{equation}
	\{ \gamma^\mu, \gamma^\nu \} \equiv \gamma^\mu \gamma^\nu + \gamma^\nu
	\gamma^\mu = 2 \eta^{\mu \nu},
\end{equation}
where $\eta^{\mu \nu}$ is the Minkowski metric. Consequently, $(\gamma^0)^2 = 1$
and $(\gamma^k)^2 = -1$, where Greek indices $\mu = 0, \dots, 3$ denote
spacetime coordinates and Latin indices $k = 1, \dots, 3$ denote spatial
components.

The chiral structure is determined by the pseudoscalar $\gamma^5 \equiv i
\gamma^0 \gamma^1 \gamma^2 \gamma^3$, where $i$ is the imaginary unit. Since
the spacetime dimension is even, $\gamma^5$ anticommutes with all vector
generators $\gamma^\mu$ and squares to the identity. This permits the
construction of the orthogonal chiral projectors
\begin{equation}
	P_L = \frac{1}{2}(1 - \gamma^5), \quad P_R = \frac{1}{2} (1 + \gamma^5),
\end{equation}
satisfying $P_L + P_R = 1$ and $P_L P_R = 0$. A spinor field $\psi$ therefore
decomposes as $\psi = P_L \psi + P_R \psi \equiv \psi_L + \psi_R$.

\subsection{Hamiltonian Formulation of the Chiral Dirac Equation}

Following the formulation of Watson and Musielak \cite{watson2020a,
watson2021a, Watson2022}, we consider the Chiral Dirac Equation (CDE) for a
spinor field $\psi$:
\begin{equation}
	(i \gamma^\mu \partial_\mu - m e^{- i \alpha \gamma^5}) \psi = 0.
	\label{eq:CDE}
\end{equation}
The corresponding effective Lagrangian density is given by
\begin{equation}
	\mathcal{L}_{\text{eff}} = \bar{\psi} (i \gamma^\mu \partial_\mu - m
	e^{- i \alpha \gamma^5}) \psi,
\end{equation}
where $\bar{\psi} = \psi^\dagger \gamma^0$ is the Dirac adjoint. We derive the
Hamiltonian density $\mathcal{H}$ via the Legendre transform, using the
canonical momentum $\pi = \partial \mathcal{L}_{\text{eff}} / \partial
\dot{\psi}$:
\begin{equation}
	\mathcal{H} = \pi \dot{\psi} - \mathcal{L}_{\text{eff}}.
\end{equation}
Explicit calculation yields:
\begin{equation}
	\mathcal{H} = \psi^\dagger (- i \gamma^{0k} \partial_k + \gamma^0 m e^{-i \alpha
	\gamma^5}) \psi.
\end{equation}
